% -----------------------------------------------------------
\bigskip
\S~\textbf{4.}
\textbf{推广到代数整数系数。}

\smallskip

\noindent
正如引言末尾已经指出的那样，如果以一个（有限）代数数域 $K$ 的代数整数代替有理整数，并以该数域的素理想 $\mathfrak{p}$ 代替素数 $p$，那么所有论证经少量修改后仍然成立。其中，\S\,1 完全不变，导出 \S\,2 中定理表述的论证也同样不变；在证明（\S\,3, 2.）以及 \S\,2, 3. 的推论中，则需要作若干补充。

\medskip\noindent\textbf{1.}
在 \S\,3, 2. 中，还须把整除多项式 $G_1(u),\ldots,G_{r-t}(u)$ 的有限多个素理想列为例外素理想（见 \S\,3, 2. 的式 (1)），因为此时不能再假定这些多项式是本原的。这样便保证了多项式 $G^*_1(u),\ldots,G^*_{r-t}(u)$（\S\,3, 2., (2)）不消失；这构成了条件 III 的加强。条件 II 保持不变，因为关于绝对不可约性的定理没有改变；条件 I 自然也同样保持不变。

\medskip\noindent\textbf{2.}
作为最终结果的 \S\,2, 3. 的推论，应按照该处的最终表述写成：

\noindent\textbf{推论。}
\textit{若 $K$ 中的代数整数 $\alpha$ 不被有限多个例外素理想中的任何一个整除，则只要 $\alpha \cdot g(x)$ 属于 $\mathfrak{I}$ 且 $g(x)$ 属于 $\mathfrak{Z}$，总可推出 $g(x)$ 属于 $\mathfrak{I}$。}

若从 $K$ 中选取代数整数 $\lambda_1,\ldots,\lambda_s$，使每个 $\lambda_i$ 均被一个例外素理想整除，但不被该素理想的平方整除，也不被其余例外素理想整除，那么这些 $\lambda$ 的幂乘积足以充当分母 $\alpha$。

\noindent\textit{证明。}
设 $(\alpha) = \mathfrak{p}_1^{\varrho_1} \cdots \mathfrak{p}_g^{\varrho_g}$ 是 $\alpha$ 的素理想分解，并且 $\mathfrak{p}_1,\ldots,\mathfrak{p}_g$ 不同于例外素理想，因而对每个 $\mathfrak{p}_i$，$\mathfrak{I}_{\mathfrak{p}_i}$ 都同构于 $\mathfrak{o}_{\mathfrak{p}_i}/\mathfrak{a}_{\mathfrak{p}_i}$。从 $K$ 的全体代数整数所成的环出发，通过在分母中加入所有且仅加入与 $(\alpha)$ 互素的代数整数，也就是不被素理想 $\mathfrak{p}_1,\ldots,\mathfrak{p}_g$ 中任何一个整除的代数整数，过渡到由 $(\alpha)$ 所定义的分式环 $\mathfrak{R}$。在 $\mathfrak{R}$ 中，素理想 $\mathfrak{p}_1,\ldots,\mathfrak{p}_g$ 仍是素理想，但变为主理想；其余所有素理想都变为单位理想；$(\alpha)$ 的分解保持不变\textsuperscript{11)}。因此，当过渡到以 $\mathfrak{R}$ 为系数环时，该推论及其证明仍完全保持 \S\,2, 3. 中的形式。

这就是说：若 $\alpha \cdot g(x)$ 属于 $\mathfrak{I}$，则 $g(x) = F(f_i(x))$，其中多项式 $F(u)$ 的系数属于 $\mathfrak{R}$。若 $\delta$ 是这些系数的一个公分母 --- 即 $\delta$ 与 $\alpha$ 互素 ---，则也可表述为：当 $\alpha \cdot g(x)$ 属于 $\mathfrak{I}$ 时，$\delta \cdot g(x)$ 属于 $\mathfrak{I}$，其中 $\delta$ 与 $\alpha$ 互素。因此 $(\delta, \alpha)$ 等于单位理想，从而得到：若 $\alpha \cdot g(x)$ 属于 $\mathfrak{I}$，则 $g(x)$ 属于 $\mathfrak{I}$。这就证明了上述推论的第一个断言。

为证明第二个断言，令 $\mathfrak{R}^*$ 为 $K$ 中由有限多个例外素理想定义的分式环，并令 $\lambda_1,\ldots,\lambda_s$ 为这些素理想在 $\mathfrak{R}^*$ 中的生成元；可以假定它们是代数整数。在 $\mathfrak{R}^*$ 中，除去 $\mathfrak{R}^*$ 的一个单位因子后，$\alpha$ 是这些 $\lambda$ 的幂乘积。因此，退回到代数整数便得到 $\gamma\alpha = \beta \lambda_1^{\varrho_1} \cdots \lambda_s^{\varrho_s}$，其中 $\gamma$ 和 $\beta$ 都不被任何例外素理想整除。由 $\alpha \cdot g(x)$ 属于 $\mathfrak{I}$，于是可推出 $\beta \cdot \lambda_1^{\varrho_1} \cdots \lambda_s^{\varrho_s} g(x)$ 属于 $\mathfrak{I}$，继而根据推论的第一个断言，还可推出 $\lambda_1^{\varrho_1} \cdots \lambda_s^{\varrho_s} g(x)$ 属于 $\mathfrak{I}$。

\begin{center}
\rule{4cm}{0.4pt}\\[0.6em]
\textbf{注释。}
\end{center}

{\small
\setlength{\parindent}{0pt}
\noindent\textsuperscript{1)} D. Hilbert，《数学问题》（Gött. Nachr. 1900，第 253--297 页），问题 14。

\noindent\textsuperscript{2)} E. Hecke，《用二元模函数构造相对阿贝尔数域》（Math. Ann. 74 (1913)，第 465--510 页），\S\,2。在 Hecke 的工作中，有限整环 $\mathfrak{I}$ 由两个变量的幂级数的三个商生成，它们之间存在一个代数关系。

\noindent\textsuperscript{3)} 参见 Hilbert 上引文献，或 E. Noether 的更详细论述《二元形式的整数不变量系统的有限性》（Gött. Nachr. 1919，第 138--156 页），\S\,1, IV。

\noindent\textsuperscript{4)} 因此，例如在一个基本形式系统的整数（射影）不变量的情形，已经证明它们可以用通常意义下的一个完全不变量系统来表示，而分母中只出现\emph{有限多个}不同的素数 --- 这一结果不能由通常的 $\Omega$-过程方法推出 ---；但这并没有对整数情形本身的有限性作出任何断言，后者迄今只在二元基本形式的情形中为人所知。（参见注\textsuperscript{3)}所引札记。）

\noindent\textsuperscript{5)} 众所周知，一个系统的“超越次数”是指这样一组代数无关函数的数目：该系统代数地依赖于它们；原文把由代数无关函数组成的系统称为“不可约系统”，本文译作“代数无关系统”。一个素理想的超越次数（维数）是指其剩余类域的超越次数。一个素理想不存在与它具有相同超越次数的真素因子（例如参见 B.\,L.\,v.\,d. Waerden，《论多项式理想的零点理论》，Math. Ann. 96 (1926)，第 183--208 页）。

\noindent\textsuperscript{6)} E. Noether，《消去理论与一般理想理论》，Math. Ann. 90 (1923)，第 229--261 页，\S\,7。借助消去理论\textsuperscript{*)}，该定理可由下述较简单的定理推出：一个绝对不可约多项式只有模有限多个素数时才会失去这一性质（A. Ostrowski，《代数量的算术理论》，Gött. Nachr. 1919，第 279--298 页，其中第 296 页。较简单的证明见 E. Noether，《绝对不可约性的一个代数判据》，Math. Ann. 85 (1922)，第 26--33 页，№ 7）。如果 $\mathfrak{I}$ 有一个整性基，其所含函数的个数恰好多于其超越次数\emph{一个} --- 例如 Hecke（注 2）所考察的情形 ---，从而 $\mathfrak{a}$ 成为主理想，那么只需利用这个较简单的定理，而无须借助消去理论。各处的素数都可以由一个（有限）代数数域的素理想代替。当 $\mathfrak{a}$ 是零理想时，该定理显然也成立。

\noindent\textsuperscript{7)} 在原来的表述中，我曾假定系数域 $P$ 是完美域；从 $P$ 过渡到 $\Omega$，把结果推广到任意域（包括非完美域），应归功于 B.\,L.\,v.\,d. Waerden。例子 $f=x^p$, $df/dx=0$ 表明，在特征为 $p$ 时，即使 $P$ 是完美域，逆命题也不成立。

\medskip
\hrule width 2.7cm
\smallskip
\noindent\textsuperscript{*)} 更确切地说，由于 $\mathfrak{a}$ 具有特征为零的系数，这一结论与该文的辅助定理 V 无关。事实上，该辅助定理的证明并不正确（在所给出的代换下，一切都可能恒等消失；正确的证明直到 B.\,L.\,v.\,d. Waerden 的《Bézout 定理的一种推广》，Math. Ann. 99 (1928)，第 497--541 页，才得以给出。参见该文注37）。

该文第 257 页还需要作一个容易看出的修正（定理 XIV 在 \S\,7 中同样没有用到）。原文说，$(\mathfrak{R})$ 中的所有元素都是通过把 $(x_i)$ 中的 $t_{\mu,\nu}$ 特殊化而得到 $(\mathfrak{R}')$ 中的元素；实际上，它们是由一个更一般的线性形式
\[
\sum s_{\lambda_1\ldots\lambda_n}(y_1)^{\lambda_1}\cdots(y_n)^{\lambda_n},
\]
经特殊化得到的，而该线性形式的指数与 $(x_i)$ 的指数相同。

\noindent\textsuperscript{8)} 关于这个雅可比行列式定理的纯代数证明 --- 这里讨论的是特征为零的情形 ---，参见 B.\,L.\,v.\,d. Waerden，《微分的代数理论》，Nieuw Archief v. Wiskunde (2) 15 (1926)，第 111--120 页。

\noindent\textsuperscript{9)} 当所讨论的是解析函数时 --- 如 Hecke 所考察的情形（注 2） ---，并且 $f_1(x),\ldots,f_t(x)$ 解析无关，而 $f_{t+1}(x),\ldots,f_r(x)$ 代数地依赖于前者时，这一附加假设得到满足。不过，这一假设被有意表述为形式性的 --- 雅可比矩阵的秩 ---，从而使它在形式定义的幂级数情形中也具有明确含义。由此可以使所有论证都保持在纯粹的形式代数范围内，而关于解析函数的定理只在特殊情形中为证明该假设得到满足时才是必要的。

\noindent\textsuperscript{10)} 例如参见 B.\,L.\,v.\,d. Waerden，《论多项式理想的零点理论》，Math. Ann. 96 (1926)，第 183--208 页，\S\,3。

\noindent\textsuperscript{11)} 例如参见 H. Grell，《代数数域和函数域中序的理论》，Math. Ann. 97 (1927)，第 524--558 页，\S\,2, 1。
}

\bigskip
\begin{center}\rule{4cm}{0.4pt}\end{center}
\begingroup
\fontspec{Noto Serif}
\begin{center}
{\Large\bfseries 论整数函数的极大整环。}\par
\vspace{0.7em}
{\bfseries 埃米·诺特（哥廷根）。}\par
\vspace{0.5em}
\emph{（摘要。）}
\end{center}

每一个没有“零因子”的环（整环）$\mathfrak{I}$ 都确定一个同类型的极大整环 $\mathfrak{M}(\mathfrak{I})$，它由所有满足下述条件的函数 $g(x)$ 组成：至少存在一个 $a\ne0$，使得 $a\cdot g(x)$ 属于 $\mathfrak{I}$。

本文在整环 $\mathfrak{I}$ 为有限生成的假设下，考察该极大整环 $\mathfrak{M}(\mathfrak{I})$ 的“分母” $a$。文中证明，总可以选择这些 $a$，使得在其全体中只有有限多个素数作为因子出现。与此同时，当 $g(x)$ 是幂级数或这类级数的商时，假定定义整环 $\mathfrak{I}$ 的函数 $f'_1(x),\ldots,f'_r(x)$ 彼此之间只能由代数关系联系。

证明的基础是把该问题归结为一般理想理论中的若干问题。这一归结通过如下方式实现：把因子 $p$ 在分母 $a$ 中的出现视为函数 $f_i$ 之间的模 $p$ 关系。

\begin{center}
(Rec. Math. XXXVI:1; 1929).
\end{center}
\endgroup

% ============================================================
\setcounter{footnote}{0}
\clearpage

